\section{Compilation setting and certificate design}
\label{sec:setting}

\subsection{Pauli representation and feasible configurations}

We work with $n$-qubit Pauli strings modulo global phase. A string is represented by a binary symplectic key $(x,z)\in\Ftwo^{2n}$; its support is the set of qubits on which $(x_q,z_q)\neq(0,0)$, and $\wt(\cdot)$ denotes support size. Two strings anticommute when their binary symplectic inner product is one.

An input $t$ contains three blocks indexed by $j\in\{0,1,2\}$, each with two Pauli targets $P_{j,0}$ and $P_{j,1}$. A feasible compilation configuration $u$ selects:
\begin{itemize}
  \item frame strings $R_{j,0}$ and $R_{j,1}$ that anticommute within every block;
  \item one shared tag string $S$;
  \item a matching, a relative assignment of the two targets to the two frames in each block, and one central-slot choice per block.
\end{itemize}
The restoration strings are $T_{j,k}=P_{j,k}R_{j,k}$, with Pauli multiplication taken modulo phase.

For three single-qubit Pauli letters, define
\begin{equation}
\phi(a,b,c)=
\begin{cases}
1, & a=b=c\neq I,\\
\one[a\neq I]+\one[b\neq I]+\one[c\neq I], & \text{otherwise}.
\end{cases}
\label{eq:phi}
\end{equation}
The unit structural support-count cost is
\begin{equation}
 c(u;t)=\sum_{j=0}^{2}\sum_{k=0}^{1}m_{j,k}\wt(R_{j,k})
 +2\wt(S)-18
 +\sum_{k=0}^{1}\sum_{q=1}^{n}\phi(T_{0,k,q},T_{1,k,q},T_{2,k,q}),
\label{eq:cost}
\end{equation}
where $m_{j,k}=2$ for a central slot and $m_{j,k}=4$ otherwise. The constant fixes the same zero point for all feasible choices and has no effect on the minimizer.

The unrestricted exact optimum and the support-at-most-two optimum are
\begin{align}
\cstar(t)&=\min_{u\in\mathcal U(t)}c(u;t),\\
\ctwo(t)&=\min_{u\in\mathcal U(t):\,\max_{j,k}\wt(R_{j,k})\leq2}c(u;t).
\end{align}
The equality $\cstar=\ctwo$ is a theorem only for Eq.~\eqref{eq:cost}; Sec.~\ref{sec:adverse-objective} gives a counterexample to transfer under reweighting.

\subsection{Forecast and explanation families}
\label{sec:certificate}

We use descriptive names for restricted feasible families. The \emph{common-anchor} family shares a simple anchor pattern across blocks. The \emph{weight-one arbitrary-anchor} family permits independent single-support anchors. The \emph{restricted-borrow} family permits a central block to borrow a restoration location subject to its original support rule. Their initial explanatory forecast is
\begin{equation}
\csimple(t)=\min\{C_{\mathrm{ca}}(t),C_{\mathrm{wa}}(t),C_{\mathrm b}(t)\}.
\label{eq:simple}
\end{equation}
Every term is a feasible-family minimum, so $\cstar\leq\csimple$; equality was initially an empirical hypothesis, not a theorem.

The exact forecast is instead
\begin{equation}
\cexact(t)=\ctwo(t).
\label{eq:exact}
\end{equation}
It is evaluated by an enumerator restricted to full frame support at most two. The forecast path does not call the unrestricted dynamic program. Exactness follows from the support-reduction theorem in Sec.~\ref{sec:support-theorem}; finite comparisons check that the implementation respects the theorem.

For interpretation, we also study an \emph{enlarged-borrow} family, which allows a borrow home outside the borrowing block's own target support, and a \emph{hybrid} family combining two tag anchors with a phantom borrow. Their costs are denoted $C_{\mathrm{b+}}$ and $C_{\mathrm h}$. The final explanation theorem is
\begin{equation}
\cstar=\ctwo=\min\{C_{\mathrm{wa}},C_{\mathrm{b+}},C_{\mathrm h}\}
\quad\text{for the unit objective.}
\label{eq:classification}
\end{equation}
These family labels explain a minimizing structure; they are not separate physical resource models.

\subsection{Chronological evidence design}

The evaluation was staged so that refutations could not be silently absorbed into later definitions:
\begin{enumerate}
  \item establish exact truth for a structured two-qubit slice and fixed application-derived rows;
  \item freeze predictions for a public-library panel and engineered boundary cases before exact evaluation;
  \item freeze a fresh seeded panel, retain the first mismatch, and bind its exact witness;
  \item introduce the theorem-backed support-two evaluator without changing the adverse row;
  \item freeze hostile shapes against the first repaired explanation;
  \item assess the hybrid family on the accumulated and adversarial finite domains;
  \item close the remaining all-size explanation sector on a complete local state domain.
\end{enumerate}

We distinguish \emph{distinct inputs} from \emph{comparison events}. The exact evaluator was compared 9,547 times on 9,546 distinct inputs because the refuting row was present in the seeded panel and was also rechecked under a separately bound audit. We report both numbers rather than treating the repeat as an independent instance.

No runtime comparison is retained. Runtime depends on the final implementation and execution environment; it is not part of the exactness claim.
